% NLP Project 2.2 - English-Italian Machine Translation
% ACM proceedings template (sigconf).
%
% 'nonacm' drops the copyright block and makes CCS/keywords optional,
% and printacmref=false removes the ACM reference-format box.
% Together these save roughly one full column - which we need for the 4-page limit.
\documentclass[sigconf,nonacm]{acmart}
\settopmatter{printacmref=false}
\renewcommand\footnotetextcopyrightpermission[1]{}
\pagestyle{plain}

\usepackage{booktabs}
\usepackage{graphicx}
\usepackage{xcolor}

% ---------------------------------------------------------------------------
% DRAFT NOTES. \DRAFTtrue prints unresolved items in red in the PDF so they
% cannot be overlooked. Switch to \DRAFTfalse before submitting: the notes then
% disappear and nothing else about the document changes.
% ---------------------------------------------------------------------------
\newif\ifDRAFT
\DRAFTfalse
\newcommand{\draftnote}[1]{\ifDRAFT\par\noindent%
  \textcolor{red}{\small\textbf{[DRAFT NOTE --- remove before submission]} #1}\par\fi}
% Look for figures both in the repo's results folders and next to this file.
% On Overleaf there is no results/ directory, so upload the PNGs into the
% project root and the second path picks them up.
% Both layouts are covered: "results/..." is the Overleaf case (main.tex at the
% project root next to results/), "../results/..." is the repo case (this file
% lives in report/). Unknown paths are simply skipped, so listing both is safe.
\graphicspath{%
  {results/task1_outputs/}{results/task2_outputs/}%
  {results/task3_outputs/}{results/task4_outputs/}{results/task5_outputs/}%
  {../results/task1_outputs/}{../results/task2_outputs/}%
  {../results/task3_outputs/}{../results/task4_outputs/}{../results/task5_outputs/}%
  {./}%
}

\begin{document}

\title{English--Italian Neural Machine Translation on Europarl}

\author{Vasileios Lalos}
\affiliation{\institution{TU Berlin}\country{}}
\author{Riccardo Tella}
\affiliation{\institution{TU Berlin}\country{}}

\begin{abstract}
We build and compare recurrent encoder--decoder models for English--Italian
translation on the Europarl v7 corpus. Starting from a vanilla RNN baseline we
show that a fixed context vector is the dominant failure mode, and address it
progressively with context feeding, LSTM units, pre-trained embeddings and
finally attention. We report BLEU and METEOR for every variant, compare both
translation directions, contrast word- and character-level models, and analyse
how performance degrades with sentence length. The strongest word model
(Italian$\rightarrow$English) reaches 9.6 BLEU, a controlled ablation shows
attention improving BLEU by 85\% over an equivalent attention-free baseline, and
an English-pivot system extends coverage to Italian$\rightarrow$Swedish, where
error propagation is the dominant limitation.
\end{abstract}

\maketitle

\section{Introduction}
% Keep this SHORT - the writing suggestions explicitly warn against long
% background sections. Three or four sentences is enough.
Machine translation maps a sentence in a source language onto a sentence in a
target language. Unlike classification there is no single correct output, so
translation quality is measured by comparing system output against reference
translations. We study English--Italian translation on the Europarl parallel
corpus, focusing on how architectural and representational choices affect
model behaviour rather than on maximising a single score.

\section{Data and Exploratory Analysis}
\label{sec:data}

We use the Europarl v7 English--Italian parallel corpus, which contains
1,909,115 line-aligned sentence pairs. Table~\ref{tab:corpus} summarises it.

\begin{table}[t]
\caption{Europarl v7 English--Italian corpus statistics (full corpus).}
\label{tab:corpus}
\begin{tabular}{lrr}
\toprule
 & \textbf{English} & \textbf{Italian} \\
\midrule
Sentence pairs      & 1,909,115 & 1,909,115 \\
Total tokens        & 49.7\,M   & 48.0\,M   \\
Vocabulary (types)  & 88,040    & \textbf{170,625} \\
Mean length (tokens)& 26.05     & 25.13 \\
Median length       & 23        & 23 \\
Max length          & 668       & 558 \\
Empty lines         & 3,257     & 4,985 \\
\bottomrule
\end{tabular}
\end{table}

The most informative asymmetry is lexical rather than positional. Italian uses
slightly \emph{fewer} tokens than English on average (25.13 vs.\ 26.05) yet
requires 94\% more distinct word types (170,625 vs.\ 88,040). This reflects
Italian's richer inflectional morphology: a single verb lemma surfaces as many
conjugated forms, each a separate entry for a word-level model, whereas English
marks the same distinctions with separate function words that are already in the
vocabulary.

This asymmetry is not merely descriptive; it shapes what a word-level model can
cover. Under the modelling pipeline used in Sections~\ref{sec:nmt}--\ref{sec:attn}
(a 20{,}000-type vocabulary over 40-token pairs), 2.50\% of Italian test tokens
fall outside the vocabulary against only 0.73\% of English ones. Both rates are
low, but the asymmetry persists: the morphologically richer target is
consistently harder to cover, and every out-of-vocabulary token is emitted as
\texttt{<unk>} and is unrecoverable, so the target side is intrinsically harder
to generate than the source side. Two predictions follow, both examined in
Section~\ref{sec:nmt}: translating \emph{into} English should be easier than into
Italian, and a character-level model, which has no word-vocabulary cap, sidesteps
this out-of-vocabulary problem entirely.

The length distributions themselves are close to identical
(Figure~\ref{fig:length}): both peak around 20 tokens with a long right tail, and
per-sentence lengths are strongly correlated across the two languages
($r = 0.946$). Length is therefore a clean axis along which to analyse errors,
which we use in Section~\ref{sec:length}. The long tail also motivates the length
cap applied in Section~\ref{sec:nmt}: a small number of very long sentences
(up to 668 tokens) would otherwise dominate padding and training cost.

\begin{figure}[t]
\centering
\includegraphics[width=\columnwidth]{length_distribution.png}
\caption{Sentence-length distributions, English vs.\ Italian (clipped at 80
tokens). The two languages are closely matched in tokens per sentence.}
\label{fig:length}
\end{figure}

Because the full corpus is too large to train on, we randomly sample 10\% of
the aligned pairs (190,912 pairs, seed 42). Sampling is performed over
\emph{indices} so that pair alignment is preserved, and the drawn indices are
stored so the sample is exactly reproducible.

\section{Preprocessing}
\label{sec:prep}

Preprocessing removes a sentence pair whenever \emph{either} side is unusable,
which preserves alignment. We lowercase, normalise repeated whitespace, strip
leading and trailing whitespace, drop empty pairs, and drop pairs containing
XML-like lines. Of 190,912 pairs, 802 (0.42\%) were removed, all of them empty;
this release of Europarl contains no XML-tag lines, so that filter was a
no-op safety check.

We deliberately \emph{retain} punctuation, numbers and stopwords, and apply
neither stemming nor lemmatisation. This is the opposite of the choice one
makes for text classification, and for a specific reason: in translation every
token must be reproduced in the output. Function words carry grammatical
information the decoder needs to generate, and stemming would destroy exactly
the inflectional endings that Italian generation depends on.

\section{Neural Machine Translation}
\label{sec:nmt}

We build recurrent encoder--decoder models that read a source sentence and
generate the target one token at a time. The encoder embeds the source tokens and
runs a recurrent network over them; its final state is passed to the decoder,
which is fed the same context at each step and predicts the next target token.
Training uses teacher forcing, cross-entropy loss that ignores padding, the Adam
optimiser and gradient clipping; decoding at test time is greedy. Every model is
evaluated with corpus BLEU and mean sentence-level METEOR (both on a 0--100
scale).

Sentence pairs longer than 40 tokens on either side are discarded, leaving
157,519 pairs, split 70/10/20 into training (110,265), validation (15,751) and
test (31,503) sets with seed 42. Vocabularies of 20,000 types per language are
built on the training split only, using a tokenizer that separates punctuation.

\subsection{Architecture and hyperparameter tuning}

We first confirmed that the recurrent unit is decisive. A vanilla RNN
encoder--decoder trained poorly and produced generic output, whereas replacing
the RNN cell with an LSTM improved every metric substantially at equal capacity.
We then tuned the embedding and hidden sizes on the validation set, selecting the
larger LSTM (embedding 128, hidden 512), and compared random against Word2Vec
initialisation. Table~\ref{tab:tuning} tracks this progression on the validation
set. The final configuration is a single-layer context-fed LSTM with embedding
dimension 128, hidden dimension 512, learning rate 0.001, batch size 128 and
Word2Vec-initialised embeddings, trained with early stopping on validation loss.

\begin{table}[t]
\caption{Model development on the validation set (English$\rightarrow$Italian).
Best per column in bold.}
\label{tab:tuning}
\begin{tabular}{lrrr}
\toprule
\textbf{Model} & \textbf{Val.\ loss} & \textbf{BLEU} & \textbf{METEOR} \\
\midrule
RNN (emb 64, hid 128)            & 4.273 & 1.89 & 13.19 \\
LSTM (emb 64, hid 128)           & 3.814 & 4.83 & 22.26 \\
LSTM (emb 128, hid 512)          & 3.531 & 6.80 & 26.68 \\
LSTM (emb 128, hid 512) + Word2Vec & \textbf{3.251} & \textbf{7.20} & \textbf{28.25} \\
\bottomrule
\end{tabular}
\end{table}

\subsection{Embeddings}
Initialising the encoder and decoder embeddings with Word2Vec (CBOW) vectors
trained on the training split, rather than learning them from scratch, improves
validation BLEU from 6.80 to 7.20 and METEOR from 26.68 to 28.25
(Table~\ref{tab:tuning}), and lowers validation loss. We therefore use Word2Vec
initialisation for the final models.

\subsection{Final results, translation direction and representation}
\label{sec:direction}
\label{sec:char}

Table~\ref{tab:nmt} reports the selected models on the held-out test set.
Translating \emph{into} English clearly outperforms translating into Italian
(BLEU 9.60 vs.\ 7.26; METEOR 33.94 vs.\ 28.12). This confirms the prediction from
Section~\ref{sec:data}: Italian is morphologically richer and carries the higher
out-of-vocabulary burden, so generating it is the harder direction.

A character-level model (same LSTM, trained on characters) removes the word
vocabulary entirely, but does not translate this into higher quality: on
English$\rightarrow$Italian it reaches only BLEU 3.86 against the word model's
7.26. Operating over much longer character sequences makes the sequence-modelling
problem harder, and the fixed context vector is a tighter bottleneck per unit of
meaning; the removal of out-of-vocabulary tokens does not compensate.

\begin{table}[t]
\caption{Final models on the test set. Best per column in bold.}
\label{tab:nmt}
\begin{tabular}{llrr}
\toprule
\textbf{System} & \textbf{Unit} & \textbf{BLEU} & \textbf{METEOR} \\
\midrule
English$\rightarrow$Italian & word      & 7.26 & 28.12 \\
Italian$\rightarrow$English & word      & \textbf{9.60} & \textbf{33.94} \\
English$\rightarrow$Italian & character & 3.86 & 19.83 \\
\bottomrule
\end{tabular}
\end{table}

\subsection{Effect of sentence length}
\label{sec:length}
Translation quality falls sharply with source length
(Figure~\ref{fig:t3length}). For the English$\rightarrow$Italian word model, BLEU
drops from 36.8 on sentences of 1--5 tokens to 8.1 on sentences of 16--20 tokens.
The character model degrades even faster (25.7 to 3.8 over the same range),
consistent with a single fixed context vector struggling to carry longer inputs.
This degradation is exactly what the attention mechanism of
Section~\ref{sec:attn} is designed to address.

\begin{figure}[t]
\centering
\includegraphics[width=\columnwidth]{task3_length.png}
\caption{Test BLEU by source-sentence length for the word and character models.
Quality declines steeply as sentences get longer.}
\label{fig:t3length}
\end{figure}

\section{Attention}
\label{sec:attn}

The models of Section~\ref{sec:nmt} compress the entire source sentence into a
single fixed vector. We replace this with \textbf{Bahdanau (additive)
attention}: at every decoding step $t$ the decoder scores all encoder states
$h_1\dots h_S$ against its current state $s_t$,
$e_{t,i} = v^\top\tanh(W_s s_t + W_h h_i)$, normalises the scores with a softmax
into weights $\alpha_{t,i}$, and consumes the weighted sum
$c_t=\sum_i \alpha_{t,i}h_i$. Scores at padded positions are set to $-\infty$
before the softmax, so padding is never attended to. We chose the additive form
over a dot product because it does not constrain encoder and decoder hidden
sizes to match, and because its weights are sharper and therefore easier to
interpret.

This is a \emph{controlled ablation}: we train a no-attention baseline and an
attention model with identical data splits, vocabularies, hyperparameters
(embedding 64, hidden 256, learning rate 0.001, batch 128, 12 epochs) and random
seeds, so their difference isolates the contribution of attention. The
configuration is a simplified, Task-3-style LSTM (20-token cap, 10k vocabulary,
whitespace tokenisation) and is \emph{not} an exact reproduction of the final
Task 3 pipeline of Section~\ref{sec:nmt} (40-token cap, 20k vocabulary, regex
tokenisation, Word2Vec, larger LSTM); the absolute scores below are therefore
not directly comparable with the Section~\ref{sec:nmt} numbers, and we do not
place them in the same table. Both systems are evaluated on the same held-out
test set.

\draftnote{Before comparing the Section~\ref{sec:nmt} and Section~\ref{sec:attn}
numbers directly, one issue remains. \textbf{BLEU scale is NOT a problem}:
Section~\ref{sec:attn} uses NLTK smoothing (\texttt{method1}) and
Section~\ref{sec:nmt} does not, but on the full test set the two are identical
(verified: 4.202 either way), because smoothing only affects n-gram orders with
zero corpus-wide matches and the converged models always produce some 4-grams.
The remaining issue is \textbf{training data}: the Section~\ref{sec:nmt} results
were produced on Colab from a stale \texttt{data/processed} that still contained
the CSV-quoting corruption (37\% of the committed Task~3 test lines are
quote-wrapped), whereas Section~\ref{sec:attn} used the repaired data. Task~3
should be re-run on the current (clean) committed data --- the notebook already
reads it, so this only needs a re-run, ideally on GPU (\texttt{DEVICE} is
currently hard-coded to \texttt{cpu}). Until then, do not put the two sections'
BLEU in one table; if it cannot be redone in time, state in the text that the
Task~3 numbers are on pre-fix data.}

\begin{table}[t]
\caption{Effect of attention, test set. Best per column in bold.}
\label{tab:attention}
\begin{tabular}{lrrrr}
\toprule
\textbf{Model} & \textbf{Val.\ loss} & \textbf{BLEU} & \textbf{METEOR} & \textbf{min} \\
\midrule
LSTM, no attention & 3.79 & 4.16 & 18.76 & \textbf{2.2} \\
LSTM + attention   & \textbf{3.36} & \textbf{7.68} & \textbf{27.07} & 5.8 \\
\bottomrule
\end{tabular}
\end{table}

Attention improves BLEU from 4.16 to 7.68 ($+85\%$ relative) and METEOR from
18.76 to 27.07 (Table~\ref{tab:attention}). The cost is training time: because
the context now depends on the decoder state, the target sequence can no longer
be pushed through the LSTM in a single call and must be decoded step by step,
making an epoch roughly $2.6\times$ slower.

\subsection{Where the gain comes from}
\label{sec:attnlength}

If the fixed context vector really is the bottleneck, its cost should grow with
sentence length, and attention should therefore help long sentences most.
Figure~\ref{fig:bleulen} confirms this prediction. The absolute gain rises
monotonically across length buckets ($+2.9$, $+2.9$, $+3.6$, $+3.8$ BLEU), and
in relative terms the effect is much starker: $+13\%$ for sentences of 1--4
tokens against $+148\%$ for sentences of 15--20. Attention does not simply make
the model uniformly better; it specifically repairs the degradation that the
single-vector bottleneck causes on longer input.

\begin{figure}[t]
\centering
\includegraphics[width=\columnwidth]{bleu_by_length.png}
\caption{BLEU by source-sentence length, with and without attention. The
advantage of attention grows with length, as the bottleneck explanation
predicts.}
\label{fig:bleulen}
\end{figure}

\subsection{What the model attends to}

Figure~\ref{fig:heatmap} shows the weights $\alpha_{t,i}$ for one test sentence.
A clear diagonal emerges, with individual weights reaching $0.9$: the model has
learned an approximately monotonic word alignment between English and Italian
even though it was never given alignment supervision. Alignment is a by-product
of learning to translate. The example also shows that alignment quality and
lexical quality are separable: the attention is well formed while the output
still contains a repetition loop (\emph{parlare di parlare di}), a common
failure of greedy decoding in an under-trained seq2seq model.

\begin{figure}[t]
\centering
\includegraphics[width=\columnwidth]{attention_heatmap.png}
\caption{Attention weights for a test sentence. Rows are generated Italian
tokens, columns are source English tokens.}
\label{fig:heatmap}
\end{figure}

Two caveats. First, the attention model reaches its best validation loss at
epoch 8 and overfits thereafter, whereas the baseline's validation loss is still
decreasing at epoch 12; both systems are reported at their best checkpoint, but
the baseline had not fully converged under the shared epoch budget. The gap
between 4.16 and 7.68 BLEU is far too large to be attributed to this. Second,
all decoding is greedy; beam search would likely raise both systems and might
narrow the gap.

\section{Pivot Translation}
\label{sec:pivot}

A pivot language is an intermediary: to translate from A to B one translates
A$\rightarrow$P and then P$\rightarrow$B. We build an Italian$\rightarrow$Swedish
system pivoting through English.

\subsection{Constructing an evaluation corpus}

Europarl distributes bilingual corpora paired with English, so Italian--Swedish
data does not exist, and \texttt{it-en} and \texttt{sv-en} are aligned
\emph{independently}: line $n$ of one is not the same sentence as line $n$ of the
other. Italian--Swedish references therefore have to be constructed before the
task can be measured at all. Because both corpora transcribe the same
parliamentary proceedings, a large share of English sentences occurs verbatim in
both, so we join the two on their normalised English side. This recovers
156,768 Italian--English--Swedish triples, of which 59,524 have all three sides
within the 20-token cap used throughout, split 70/10/20 into 41,668 training,
5,952 validation and 11,904 test triples. When two Italian sentences share the
same English side we keep the first Swedish match; about 4.5\% of the resulting
triples are exact duplicates (repeated boilerplate in the proceedings), which we
did not remove, so a small share may be shared across the train/test split. This
mildly inflates the absolute scores but does not affect the comparisons below,
which are all made on the same split.

\subsection{Systems}

All four systems use the attention architecture of Section~\ref{sec:attn} with
identical hyperparameters, initialisation and splits, and each is evaluated on
the same held-out triples. Besides the two legs and the pivot we train a
\emph{direct} Italian$\rightarrow$Swedish model, and we additionally measure the
second leg on \emph{reference} English rather than on the first leg's output.
That last configuration is an oracle: it shares the second-leg model with the
pivot and differs only in the quality of its input, so the difference between
them isolates the error inherited from the first leg. The pivot itself is fed the
first leg's own output, since using reference English would leak information no
pivot has at inference time.

\begin{table}[t]
\caption{Pivot translation, test set. The oracle row shares its model with the
pivot row and differs only in whether the English input is the reference or the
first leg's output.}
\label{tab:pivot}
\begin{tabular}{llrr}
\toprule
\textbf{System} & \textbf{Target} & \textbf{BLEU} & \textbf{METEOR} \\
\midrule
IT$\rightarrow$EN (leg 1)              & EN & 10.68 & 32.23 \\
EN$\rightarrow$SV (leg 2, oracle input)& SV & \textbf{11.00} & \textbf{35.12} \\
IT$\rightarrow$EN$\rightarrow$SV (pivot)& SV & 5.95 & 22.86 \\
IT$\rightarrow$SV (direct)             & SV & 6.06 & 25.08 \\
\bottomrule
\end{tabular}
\end{table}

\subsection{Error propagation}

Taken separately the two legs are comparable in quality (BLEU 10.68 and 11.00),
but chaining them yields only 5.95 --- the pivot retains just \textbf{54\% of the
second leg's oracle performance}. The 5.05 BLEU lost between the oracle and the
pivot cannot be attributed to the second-leg model, which is identical in both
rows; it is caused entirely by errors the first leg makes and the second leg then
treats as correct input. Figure~\ref{fig:pivot} shows the decomposition. This is
the characteristic weakness of pivoting: quality is roughly multiplicative along
the chain rather than additive, so two individually usable systems compose into a
markedly weaker one, and no amount of second-leg capacity can recover information
already destroyed upstream.

\begin{figure}[t]
\centering
\includegraphics[width=\columnwidth]{pivot_error_propagation.png}
\caption{Decomposing the pivot. The first bar is the second leg given perfect
English; the second is the same model given the first leg's output. The gap is
inherited error.}
\label{fig:pivot}
\end{figure}

\subsection{Pivot versus direct translation}

On BLEU the pivot is indistinguishable from the direct system (5.95 vs.\ 6.06, a
gap of 0.11), which is a stronger result for pivoting than we expected. METEOR
disagrees, favouring the direct system by 2.22 points (22.86 vs.\ 25.08). The
disagreement is informative rather than contradictory: BLEU rewards exact n-gram
matches and both systems are weak enough that few 4-grams match either way,
whereas METEOR credits stem-level unigram overlap and so detects that the direct
system recovers more of the reference content. Where the metrics do agree, the
direct system is at least as good, which is the expected ordering.

Two caveats qualify the comparison. First, the direct system is not a strong
ceiling: it is trained on the same 41,668 triples as the legs, so it is limited
by data volume rather than by the absence of an intermediate language, and a
direct system trained on a larger Italian--Swedish corpus would very likely
extend its advantage. Second, and for the same reason, the two legs here are
trained on Italian--Swedish-aligned data, whereas the practical motivation for
pivoting is precisely the case where no such data exists; a deployment would
train each leg on its full bilingual corpus. Both choices were made to keep every
system on an identical split so the comparison isolates the pivot mechanism, and
they should be read as bounding the experiment rather than the method.

The practical trade-off is nevertheless clear. Pivoting costs two models and two
decoding passes at inference, and roughly halves quality relative to a
well-supplied second leg, but it reaches a language pair for which no direct
training data exists at all --- and here it did so at essentially no BLEU cost
relative to what direct training on the available data achieved.

\section{Conclusion}

Across all experiments the recurrent unit and the fixed context vector were the
decisive factors. Replacing the vanilla RNN with an LSTM and adding Word2Vec
initialisation gave the largest gains among the word models, and attention, which
removes the single-context bottleneck, produced the clearest improvement of all,
most strongly on long sentences. The data itself set the harder direction:
Italian's richer morphology makes it the more difficult target, so
Italian$\rightarrow$English consistently outperformed the reverse. Pivoting
through English reached a language pair with no direct data, at the cost of
propagated first-leg errors. Natural next steps are attention combined with the
full 40-token pipeline, sub-word tokenisation to cut the remaining
out-of-vocabulary loss, and beam search in place of greedy decoding.

\end{document}
